\documentclass[12pt,a4paper]{article}

% --- LaTeX Packages ---
\usepackage[UTF8,scheme=chinese]{ctex}
\usepackage{geometry}
\usepackage{amsmath,amssymb,amsfonts}
\usepackage{booktabs}
\usepackage{float}
\usepackage{listings}
\usepackage{color}
\usepackage{fancyhdr}
\usepackage{array}
\usepackage{titlesec}
\usepackage{hyperref}

% A4 paper, margins: top/bottom/left/right 2.5cm
\geometry{top=2.5cm,bottom=2.5cm,left=2.5cm,right=2.5cm}

% Chinese font settings
\setCJKmainfont{SimSun}[BoldFont=SimHei, ItalicFont=KaiTi]
\setCJKsansfont{SimHei}

% Line spacing
\renewcommand{\baselinestretch}{1.3}

% Section spacing and styling
\titleformat{\section}{\zihao{-3}\bfseries\heiti}{\thesection}{1em}{}
\titleformat{\subsection}{\zihao{4}\bfseries\heiti}{\thesubsection}{1em}{}

% Listings configuration
\definecolor{dkgreen}{rgb}{0,0.6,0}
\definecolor{gray}{rgb}{0.5,0.5,0.5}
\lstset{
  frame=tb,
  language=Python,
  aboveskip=3mm,
  belowskip=3mm,
  showstringspaces=false,
  columns=flexible,
  basicstyle={\linespread{1.0}\footnotesize\ttfamily},
  keywordstyle=\color{blue},
  commentstyle=\color{dkgreen},
  breaklines=true,
  tabsize=4
}

\begin{document}

\begin{center}
    \vspace*{0.5cm}
    {\zihao{2}\bfseries\heiti 全国大学生数学建模竞赛}\\
    {\zihao{3}\bfseries\heiti 人工智能工具使用详情说明书}\\[0.6cm]
    \rule{\textwidth}{1.5pt}\\[0.8cm]
\end{center}

\section{所用 AI 工具基本信息}
根据《全国大学生数学建模竞赛人工智能工具使用规定（2025年试行）》的要求，本参赛队在竞赛期间规范、公开、透明地使用了人工智能工具。相关工具的基本信息如表\ref{tab:ai_info}所示：

\begin{table}[H]
    \centering
    \caption{参赛队所使用的人工智能（AI）工具基本信息表}
    \label{tab:ai_info}
    \begin{tabular}{cccc}
        \toprule
        \textbf{工具名称} & \textbf{版本/型号} & \textbf{开发机构/公司} & \textbf{使用日期} \\ \midrule
        DeepSeek & deepseekV4-pro & 深度求索 (DeepSeek) & 2026年5月24日 \\
        \bottomrule
    \end{tabular}
\end{table}

\section{具体使用目的与环节}
本参赛队仅在学术思路排查、部分数学公式的物理背景探讨、LaTeX格式的排版美化以及求解程序的辅助语法调试中使用了 AI 工具。所有核心模型的构建、核心公式的离散化推导、程序代码的算法逻辑实现与图表生成，均由参赛队全权独立完成。具体使用环节如表\ref{tab:ai_usage}所示：

\begin{table}[H]
    \centering
    \caption{AI 工具的具体使用目的与涉及环节}
    \label{tab:ai_usage}
    \begin{tabular}{lp{6.5cm}p{6cm}}
        \toprule
        \textbf{涉及环节} & \textbf{使用 AI 工具的具体目的} & \textbf{采纳与人工修改情况说明} \\ \midrule
        数学公式理顺 & 探讨等效热参数 (ETP) 模型的连续微分方程在大步长离散化时的潜在发散风险。 & 采纳了 AI 提出的使用指数积分法代替经典 Euler 离散的思路。由队员手写推导了具体的解析形式。 \\ \addlinespace
        代码语法调试 & 辅助排查 Python 中 PuLP 线性规划求解器在添加大规模绝对值约束时的语法报错。 & 采纳了 AI 关于引入辅助决策变量进行非线性约束线性化的建议。队员独立编写并运行了完整的求解代码。 \\ \addlinespace
        论文格式理顺 & 美化 LaTeX 论文中 TikZ 系统流程图的连线排版和 listings 代码段的单倍行距设置。 & 采纳了 AI 给出的 TikZ styling 建议代码。队员根据竞赛排版规范进行手写参数微调和两轮编译。 \\
        \bottomrule
    \end{tabular}
\end{table}

\clearpage

\section{关键交互记录}
以下展示参赛队与人工智能工具交互中的两个关键技术记录，用以说明工具在启发思路和排查语法中的辅助作用。

\subsection{关键交互记录 1：关于 ETP 模型大步长离散化的发散排查}
\begin{description}
    \item[参赛队提示词 (Prompt)：] \ \\
    “我们在建立空调负荷等效热参数 (ETP) 模型时，一阶微分方程为：
    \[ C \frac{dT_{in}(t)}{dt} = \frac{T_{out}(t) - T_{in}(t)}{1000 R} - EER \cdot P(t) \]
    当我们将仿真步长设为 $\Delta t = 1\ \text{min}$ 时，采用向前欧拉离散法在仿真中出现了温度数值发散和大幅震荡。这在物理上是不合理的。请问有什么方法可以从根本上解决大步长下的离散化数值不稳定性？”
    
    \item[AI 工具回复 (Response)：] \ \\
    “向前欧拉法是显式格式，其数值稳定性要求步长必须极小。对于时间常数较小或热阻较小的建筑 ETP 系统，$\Delta t = 60\ \text{s}$ 会超出其稳定域，导致强烈的数值震荡。\\
    **根本解决方案**是使用**解析指数积分法 (Analytical Exponential Integrator)**。由于在一步区间 $[\tau_k, \tau_{k+1}]$ 内室外温度 $T_{out}$ 和空调功率 $P$ 可近似视为常数，该线性非齐次微分方程具有解析精确解。其更新格式为：
    \[ T_{in}(k+1) = a T_{in}(k) + (1-a) T_{out}(k) - b P(k) \]
    其中：
    \[ a = e^{-\frac{\Delta t}{1000 R C}}, \quad b = 1000 R \cdot EER \cdot (1-a) \]
    此格式属于半隐式解析更新，在任何离散步长下都绝对稳定且无截断误差，能完美解决数值发散。”
    
    \item[参赛队采纳与人工修改情况：] \ \\
    参赛队完全采纳了此建议。针对 B 类建筑参数，手写推导了精确的离散化系数 $a$ 和 $b$，并编写了仿真更新代码。实践证明：该格式在 24 小时常规运行仿真中极其稳定，消除了任何震荡发散，完全符合物理实际。
\end{description}

\clearpage

\subsection{关键交互记录 2：关于 PuLP 线性规划求解器中非线性绝对值约束的处理}
\begin{description}
    \item[参赛队提示词 (Prompt)：] \ \\
    “我们在使用 Python 中的 PuLP 库编写多用户空调集群优化调度程序时，目标函数需要最小化用户的温度不适度，即室内实际温度与设定值 Tup 的差值绝对值累加之和：
    \[ \min \sum |T_g(t) - T_{up, g}| \]
    但在 PuLP 中直接写 `abs` 函数会报错，因为它是非线性的。请问如何在保持线性规划 (LP) 求解器高效性的前提下，将这个绝对值项线性化？”
    
    \item[AI 工具回复 (Response)：] \ \\
    “在数学规划中，最小化绝对值和 $\min |x|$ 是非常经典的非线性项，可以通过引入**辅助决策变量 (Auxiliary Variables)** 进行等价的线性化处理。\\
    对于每一个群组 $g$ 的每一次时间步长 $t$，定义一个非负辅助决策变量 $z_g(t) \ge 0$。将目标函数中的绝对值项替换为 $z_g(t)$。然后，对于所有的 $g$ 和 $t$，在 PuLP 中添加以下两条线性不等式约束：
    \[ z_g(t) \ge T_g(t) - T_{up, g} \]
    \[ z_g(t) \ge 0 \]
    这样，线性求解器在最小化 $z_g(t)$ 时，会迫使其自动等于 $\max(0, T_g(t) - T_{up, g})$（即单侧越限惩罚）。这种转换完全保持了约束系统的线性，因此能够利用高效的 simplex 或内点法在数秒内解出全局最优调度方案。”
    
    \item[参赛队采纳与人工修改情况：] \ \\
    参赛队深入学习并完全采纳了辅助变量的线性化机制。在 \lstinline!solve_problem_a.py! 的 PuLP 建模部分定义了辅助字典变量 \lstinline!Z_vars!，并成功写入了上述两个线性不等式约束。在实际求解中，1000 户居民在 360 分钟时域下的 36 万个决策变量的 LP 问题仅用 0.8 秒便解算完成，求得了全局最优调度。
\end{description}

\section{总结与学术诚信承诺}
本参赛队秉承诚实守信、公开透明的学术态度。在此郑重承诺：
\begin{enumerate}
    \item 本篇论文的全部核心数学模型（ETP指数离散化、群组降维聚合、大规模空调集群优化调度、机会约束鲁棒调度）均由参赛队员独立讨论、构建并完成推导。
    \item 全部的计算与绘图 Python 代码均由参赛队员自主调试运行，没有使用任何 AI 工具自动生成完整代码或支撑材料。
    \item 参赛队在使用 AI 辅助的过程中，完全遵循了数学建模竞赛组委会的试行规定，并在参考文献与支撑材料中进行了如实和规范的批注。
\end{enumerate}

\vspace{1.5cm}
\begin{flushright}
    参赛队组号：13165 \\
    参赛队员签名：严浩睿、陈俊明、吴扬 \\
    日期：2026年5月24日
\end{flushright}

\end{document}
